\section{The original mixed native square on its coordinate section}

Let $k$ be a field, let $A$ be a commutative $k$-algebra, and let
$x_1,\ldots,x_n$ define an etale coordinate map
$k[X_1,\ldots,X_n]\to A$. Fix the original finite root family:
positive degrees $N_i$, units $u_i\in A^\times$, a set $I$ of active
indices, and an injection $\sigma:I\hookrightarrow\{1,\ldots,n\}$.
Assume each $N_i$ is a unit in $A$. Let $B$ be the entire original
mixed root algebra, $S$ the entire unit-root algebra, and
$C=S\otimes_A B$. Let $P$ be the original pure root algebra over $S$,
and let $\rho:C\to P$ be the produced purification homomorphism.
All original root factors and all components of these algebras are retained.

Write $M_A,M_B,M_S,M_P$ for the original native affine top exterior
cotangent section modules, and $M_C$ for the original tensor comparison's
native section module with its original composite field structure.
Scalar restrictions use the original coefficient homomorphisms.
Let $\omega_x$ be the original sheafification of
$dx_1\wedge\cdots\wedge dx_n$ in $M_A$. For a tuple $z$ in $P$,
write $\omega_P(z)$ for its native sheafified coordinate wedge.
The notation $\operatorname{read}_P$ below forgets only the original
restriction-of-scalars presentation of the same section.

\begin{definition}[The original coordinate input and native readers]
\label{def:kn-mixed-native-coordinate-paths}
\lean{EnvelopeProperty.Vanishing.NadelGlobalCanonicalTrace.mixedNativeCanonicalSquare_inputValue,EnvelopeProperty.Vanishing.NadelGlobalCanonicalTrace.mixedNativeCanonicalSquare_restrictedInputValue,EnvelopeProperty.Vanishing.NadelGlobalCanonicalTrace.mixedNativeCanonicalSquare_purificationInputCoordinates}
\uses{lem:kn-nadel-mixed-original-native-canonical-input,lem:kn-nadel-mixed-comparison-native-tensor,lem:kn-mixed-original-native-value,lem:kn-unit-original-native-value}
\leanok
Let $M_{\mathrm{ctor}}$ and $\alpha:S\to C$ be the exact native
module and coefficient map used by the original purification constructor.
Let $v_{\mathrm{ext}}$ be its original extended coordinate-point
readback, and retain its restricted presentation
\[
 \zeta=\operatorname{fromUnderlying}_{\alpha,M_{\mathrm{ctor}}}
                   (v_{\mathrm{ext}}).
\]
Let $\mathcal R_\rho$ be the original outer-coefficient native Gamma
reader in the purification coordinate theorem, followed by its original
target readback. Its source is this exact restricted module and its
target is the original native $P$-section carrier.
Let $U:S\otimes_A M_A\simeq M_S$ be the original unit comparison,
$d_p:M_S\to\operatorname{Res}_{S\to P}M_P$ the original pure native
pullback, and $\eta_S$ the actual extension--restriction adjunction unit
on the original base section module, read in its original extension
carrier. The purification input tuple is exactly
$j\mapsto i_B(\operatorname{im}_{A\to B}(x_j))$.
\end{definition}
\begin{proof}
Read the coefficient map, native source module and coordinate value
from the original constructor. Restricting its scalar presentation keeps
that same section. The reader is the original function of this exact
restricted source, with its original field-transport witness retained.
The unit comparison and pure differential are the original native maps.
These definitions introduce no new compatible map or normalization
equation.
\end{proof}

\begin{lemma}[The mixed path on the original coordinate section]
\label{lem:kn-mixed-native-coordinate-left-value}
\lean{EnvelopeProperty.Vanishing.NadelGlobalCanonicalTrace.mixedNativeExtended_originalCoordinateValue,EnvelopeProperty.Vanishing.NadelGlobalCanonicalTrace.mixedNativeCanonicalSquare_purification_input_coordinate_typed,EnvelopeProperty.Vanishing.NadelGlobalCanonicalTrace.mixedNativeCanonicalSquare_left_coordinate}
\uses{def:kn-mixed-native-coordinate-paths,lem:kn-mixed-original-native-readback,lem:kn-global-trace-native-coordinate-wedge}
\leanok
Let $i_B:B\to C$ be the original tensor-right homomorphism and set
\[
 z^L_j=\rho\bigl(i_B(\operatorname{im}_{A\to B}(x_j))\bigr).
\]
Then the actual original reader satisfies
\[
 \mathcal R_\rho(\zeta)=\omega_P(z^L).
\]
The input tuple to the purification comparison is exactly
$j\mapsto i_B(\operatorname{im}_{A\to B}(x_j))$ in the original $C$.
\end{lemma}
\begin{proof}
The original mixed native pullback sends $\omega_x$ to the coordinate
wedge of the original images of $x_j$ in $B$. The original extension
point theorem computes its coordinate value using the actual unit and
sheaf differential, with coefficient the image of $1$. Exterior
naturality sends that wedge to the tensor-right image coordinate wedge
in the original $C$; the coefficient $1$ acts as the identity.

Read that value in the exact original input of the purification
coordinate theorem. The native/raw and restriction-of-scalars readers
retain the same section. The accepted original-source equality is
applied before the original purification reader; thus the actual
purification coordinate identity applies to this same tuple. It sends
the wedge to the wedge of its images under $\rho$, with the original
field transport retained. This is exactly the displayed value in $M_P$.
\end{proof}

\begin{lemma}[The pure path on the same coordinate section]
\label{lem:kn-mixed-native-coordinate-right-value}
\lean{EnvelopeProperty.Vanishing.NadelGlobalCanonicalTrace.mixedNativeCanonicalSquare_unit_congr_coordinate,EnvelopeProperty.Vanishing.NadelGlobalCanonicalTrace.mixedNativeCanonicalSquare_right_coordinate,EnvelopeProperty.Vanishing.NadelGlobalCanonicalTrace.mixedNativeCanonicalSquare_right_target_readback,EnvelopeProperty.Vanishing.NadelGlobalCanonicalTrace.mixedNativeCanonicalSquare_original_right_coordinate}
\uses{def:kn-mixed-native-coordinate-paths,lem:kn-unit-original-native-value,lem:kn-global-trace-native-coordinate-wedge}
\leanok
Let $i_S:S\to P$ be the original structural homomorphism and set
\[
 z^R_j=i_S(\operatorname{im}_{A\to S}(x_j)).
\]
The actual pure value through the original target reader satisfies
\[
 \operatorname{read}_P\bigl(d_p(U(\eta_S(\omega_x)))\bigr)
   =\omega_P(z^R).
\]
Both the pure target's original names and its restricted native section
presentation retain this same value in the same original $P$.
\end{lemma}
\begin{proof}
The original unit-cover comparison evaluates the actual adjunction unit
of $\omega_x$ as
the coordinate wedge of the images of the $x_j$ in $S$. Applying the
actual pure native differential gives the wedge of their original
images under $S\to P$, by the native algebra-map coordinate identity.
Restriction of scalars keeps the underlying section. Reading the same
original pure target through its two defining names leaves the value
unchanged. These original reader identities give the displayed formula.
\end{proof}

\begin{lemma}[The produced ring square identifies the actual target wedges]
\label{lem:kn-mixed-native-coordinate-target-wedge}
\lean{EnvelopeProperty.Vanishing.NadelGlobalCanonicalTrace.mixedNativeCanonicalSquare_original_target_tuple,EnvelopeProperty.Vanishing.NadelGlobalCanonicalTrace.mixedNativeCanonicalSquare_wedge_coordinate}
\uses{lem:kn-mixed-purification-original-ring-square,lem:kn-mixed-native-coordinate-left-value,lem:kn-mixed-native-coordinate-right-value}
\leanok
The two literal original tuples in $P$ agree:
$z^L=z^R$. Consequently
$\omega_P(z^L)=\omega_P(z^R)$ in the original native target module.
\end{lemma}
\begin{proof}
The original produced ring square is
\[
 \rho\circ i_B\circ(A\to B)=i_S\circ(A\to S).
\]
Evaluate it at every original $x_j$. This gives equality of the two
target tuples at every index, hence equality of those tuples in $P$.
Apply the same native coordinate-wedge construction to this equality.
No coordinate, coefficient ring or field structure is replaced in
passing from the tuple equality to the wedge equality.
\end{proof}

\begin{theorem}[The original mixed native coordinate equation]
\label{thm:kn-mixed-native-original-coordinate-equation}
\lean{EnvelopeProperty.Vanishing.NadelGlobalCanonicalTrace.mixedNativeCanonicalSquare_coordinate}
\uses{lem:kn-mixed-native-coordinate-left-value,lem:kn-mixed-native-coordinate-right-value,lem:kn-mixed-native-coordinate-target-wedge}
\leanok
For the original etale coordinates, positive tame degrees, units and
active-coordinate injection above, the two actual native paths have
the following exact original-reader coordinate equation:
\[
 \mathcal R_\rho(\zeta)
 =\operatorname{read}_P\bigl(d_p(U(\eta_S(\omega_x)))\bigr).
\]
The field may have arbitrary characteristic. No integrality,
nontriviality or connectedness hypothesis on $A$, nor a caller
coordinate-compatibility, trace or vanishing equation, is assumed.
The conclusion is the displayed coordinate-value equation.
\end{theorem}
\begin{proof}
The original-reader identity gives the left side as $\omega_P(z^L)$.
The produced ring square gives $z^L=z^R$, hence equality of their
native target wedges. The original target readback and target-name
identities identify that value with the right-path expression.
Composing these actual equalities proves the stated reader equation.
An identification with separately named whole-map expressions evaluated
on a tensor generator is a further presentation comparison. No such
comparison, equality of entire homomorphisms, or counit normalization is
asserted here.
\end{proof}
